\chapter{Хопфова линия на ориентированных переходах и условное переоткрытие}

После представительного запрета $SU(2)_F$ и запрета одной двойки
$\operatorname{Spin}^h$ казалось, что спинорный знак неизбежно требует
нового внутреннего модуля. Повторное чтение раннего текста для второй
статьи на Хабре изменило постановку. В нём первичным объектом была не
дополнительная группа, а нить-переход в полном расслоении Хопфа.

Настоящий гейт проверяет две разные гипотезы, которые нельзя смешивать:
\begin{enumerate}
 \item достаточно ли одной плоской фазы порядка четыре;
 \item содержит ли уже имеющееся накрытие $S^3\to\mathbb{RP}^3$ полную
 хопфову линию с единичным классом Черна, которую можно назначить
 ориентированной стрелке Мориты.
\end{enumerate}

\section{Почему одной фазы порядка четыре недостаточно}

Пусть внутренний цикл равен $C_4$ и
\begin{equation}
 \chi(k)=i^k,
 \qquad k\in\mathbb Z_4.
\end{equation}
Это обычный характер:
\begin{equation}
 \chi(a+b)=\chi(a)\chi(b).
\end{equation}
Поэтому построенный из него проективный множитель
\begin{equation}
 \omega(a,b)=\frac{\chi(a)\chi(b)}{\chi(a+b)}
\end{equation}
тождественно равен единице. Последовательность
$1,i,-1,-i,1$ является одномерным представлением, а не нетривиальным
двухкоциклом.

Соответствующая плоская линия имеет нулевую кривизну и
\begin{equation}
 c_1=0.
\end{equation}
Следовательно, одна дискретная фаза может хранить знак и задержку, но не
может породить хопфову линию с индексом один.

\begin{nogocriterion}[Запрет плоской фазы]
Фаза порядка четыре без неплоской связности не даёт ненулевого первого
класса Черна. Поэтому прежний четырёхтактный цикл сам по себе не закрывает
индексный гейт.
\end{nogocriterion}

\section{Полное расслоение Хопфа}

Ранний текст содержал более сильный объект:
\begin{equation}
 U(1)\longrightarrow S^3\longrightarrow S^2.
 \label{eq:v5-hopf-fibration}
\end{equation}
Пусть на северной карте $S^2$ выбран нормированный спинор
\begin{equation}
 z_N(\theta,\phi)=
 \begin{pmatrix}
  \cos(\theta/2)\\
  e^{i\phi}\sin(\theta/2)
 \end{pmatrix}.
 \label{eq:v5-hopf-north-section}
\end{equation}
Тогда
\begin{equation}
 n_a=z_N^\dagger\sigma_a z_N,
 \qquad
 P_+=z_Nz_N^\dagger=\frac{I+n_a\sigma_a}{2}.
 \label{eq:v5-hopf-projector}
\end{equation}

На южной карте можно взять
\begin{equation}
 z_S=e^{-i\phi}z_N.
\end{equation}
Переходная функция имеет единичное наматывание. Каноническая связность
Берри
\begin{equation}
 A=-iz^\dagger dz
\end{equation}
имеет кривизну
\begin{equation}
 F=\frac12\sin\theta\,d\theta\wedge d\phi,
\end{equation}
и потому
\begin{equation}
 c_1(L)=\frac1{2\pi}\int_{S^2}F=1.
 \label{eq:v5-hopf-c1-one}
\end{equation}
Это в точности та положительная собственная линия массы
$n\cdot\boldsymbol\sigma$, которая в проекторном гейте дала единичный
индекс.

\section{Где находится фаза порядка четыре}

Вертикальное действие хопфова волокна имеет вид
\begin{equation}
 z\longmapsto e^{i\alpha}z.
\end{equation}
При $\alpha=\pi/2$ получаем
\begin{equation}
 z\mapsto iz\mapsto-z\mapsto-iz\mapsto z.
\end{equation}
При этом $n_a$ и $P_+$ не меняются. Следовательно, прежняя фаза порядка
четыре не создаёт хопфово расслоение, но естественно вкладывается в его
полное $U(1)$-волокно как дискретная четверть оборота.

Это исправляет раннюю метафору: мёбиусный знак является не всей
топологией, а наблюдаемым элементом более богатого волокна.

\section{Накрытие и quotient-чтение}

Геометрическое ядро проекта уже различает
\begin{equation}
 S^3\simeq SU(2),
 \qquad
 \mathbb{RP}^3\simeq SO(3)=SU(2)/\{\pm1\}.
\end{equation}
Хопфово $S^3$ является круговым расслоением над $S^2$ с классом Черна
один. После quotient по центральному $\mathbb Z_2$ получается
$\mathbb{RP}^3=L(2,1)$, являющееся круговым расслоением с классом Черна
два. Поэтому
\begin{equation}
 \boxed{
 \text{спинорное накрытие }S^3\Rightarrow c_1=1,
 \qquad
 \text{векторный quotient }\mathbb{RP}^3\Rightarrow c_1=2.}
 \label{eq:v5-hopf-cover-quotient-c1}
\end{equation}

Это не новое случайное пространство. Оба носителя уже присутствуют в
RPFT как спинорное и векторное чтения одной геометрии. Ранее они
использовались для спектров и спин-структур; теперь обнаружено возможное
индексное чтение их различия.

\section{Линия Фелла на группоиде переходов}

Пусть $L\to S^2$ --- хопфова линия. Для стрелки $x\to y$ определим
одномерное пространство переноса
\begin{equation}
 \mathcal F_{y\leftarrow x}=L_y\otimes L_x^*.
 \label{eq:v5-hopf-fell-arrow}
\end{equation}
Для составимых стрелок имеется каноническое умножение
\begin{equation}
 \mathcal F_{z\leftarrow y}\otimes
 \mathcal F_{y\leftarrow x}
 \longrightarrow
 \mathcal F_{z\leftarrow x},
 \label{eq:v5-hopf-fell-composition}
\end{equation}
получаемое сокращением $L_y^*\otimes L_y$. Это простейшая линия Фелла над
группоидом переходов. Общая теория центральных расширений группоидов и
скрученных гильбертовых расслоений описывает именно такие
мультипликативные линии; см. path{arXiv:0711.1906},
path{arXiv:math/0306138} и path{arXiv:1801.03016}.

Важно различать два источника данных. Ассоциативность композиции не
требует свободного коэффициента, но индекс приходит из $c_1(L)=1$, а не
из плоского характера $C_4$.

\section{Совместимость с KO6 без нового семейного дублета}

Прямая стрелка соответствия Мориты имеет тип
\begin{equation}
 E:M_{15}\longrightarrow M_{20},
\end{equation}
а обратная --- $E^*$. Кандидат на новое чтение имеет вид
\begin{equation}
 E\otimes L,
 \qquad
 E^*\otimes L^*.
 \label{eq:v5-hopf-oriented-morita-pair}
\end{equation}
Комплексное сопряжение меняет
\begin{equation}
 c_1(L)=+1
 \quad\longleftrightarrow\quad
 c_1(L^*)=-1.
\end{equation}
Поэтому вещественная структура KO6 может обменивать две направленные
ветви вместе с $E\leftrightarrow E^*$. Полный вещественный модуль остаётся
парным, тогда как ориентированное физическое чтение одной стрелки видит
единичный индекс.

Здесь не вводится фиксированный внутренний дублет
$\mathbb C^2_F$ и не объявляется новая калибровочная группа $SU(2)_F$.
Локальный двухкомпонентный $z$ является координатой подъёма хопфовой
линии с $U(1)$-избыточностью, а не набором двух новых наблюдаемых
фермионов. Поэтому представительный и $\operatorname{Spin}^h$-запреты не
применяются к этому объекту.

\section{Главный оставшийся разрыв}

Математическая линия существует и уже содержится в геометрическом
инвентаре проекта. Но пока не доказано, что направление стрелки
$E/E^*$ функториально выбирает ориентированную ось $n/-n$ проекторного
дефекта. Сам факт, что $E$ и $E^*$ имеют противоположные источник и цель,
ещё не является законом
\begin{equation}
 E\longmapsto n,
 \qquad
 E^*\longmapsto-n.
 \label{eq:v5-hopf-missing-orientation-functor}
\end{equation}

Если такое отображение придётся назначить вручную после получения
индекса, новая конструкция станет переименованной подгонкой. Оно должно
следовать из градуировки, высоты, полярной декомпозиции перехода или
родительской суперсвязности до вычисления наблюдаемого индекса.

\begin{theorem}[Условное хопфово переоткрытие]
Плоская фаза порядка четыре не создаёт индексный класс. Однако полное
хопфово расслоение на уже существующем спинорном накрытии $S^3$ даёт
каноническую линию $L$ с $c_1=1$. Пара
$E\otimes L$, $E^*\otimes L^*$ совместима с комплексным сопряжением и не
требует нового фиксированного семейного дублета. Физическое переоткрытие
остаётся условным до вывода ориентационного функтора
\eqref{eq:v5-hopf-missing-orientation-functor}.
\end{theorem}

\begin{protocol}
\begin{enumerate}
 \item чистый $C_4$-характер как нетривиальный двухкоцикл:
 \textbf{отклонён};
 \item класс Черна плоской фазы: \textbf{нуль};
 \item хопфова линия на $S^3\to S^2$: \textbf{построена};
 \item первый класс Черна: \textbf{$c_1=1$};
 \item фаза порядка четыре внутри $U(1)$-волокна:
 \textbf{реализована};
 \item связь спинорного накрытия и векторного quotient:
 \textbf{$1\to2$ по классу Черна};
 \item линия Фелла на стрелках: \textbf{типизирована};
 \item KO6-пара $L/L^*$: \textbf{совместима};
 \item новый фиксированный $\mathbb C^2_F$: \textbf{не требуется};
 \item функтор $E/E^*\mapsto n/-n$: \textbf{не выведен};
 \item статус ветви: \textbf{условно переоткрыта одним гейтом};
 \item физическое замыкание: \textbf{не достигнуто}.
\end{enumerate}
\end{protocol}

Машинный протокол сохранён в
\path{s2t_v5_hopf_fell_line_transition_lift_gate_results.json}.